%
\newacronym{mat}{Matlab}{Matrix Laboratory}
\newacronym{ide}{IDE}{Integrated Development Environment}

% 1. Definisce un nuovo stile di teorema
\newtheoremstyle{solotitolo}
{} % Spazio sopra
{} % Spazio sotto
{\normalfont} % Font del testo interno (usa \itshape se preferisci il corsivo)
{} % Indentazione
{\bfseries} % Font del titolo
{} % Punteggiatura dopo il titolo (es. inserisci "." se vuoi un punto)
{ } % Spazio dopo il titolo (usa { } se vuoi il testo sulla stessa riga)
{\thmnote{#3}} % LA MAGIA: Stampa SOLO il parametro passato tra le parentesi quadre

% 2. Applica lo stile e crea l'ambiente (usa l'asterisco per togliere i contatori di default)
\theoremstyle{solotitolo}
\newtheorem*{formula}{}

% 3. Il tuo stile tcolorbox (invariato)
\tcbset{
    formulastyle/.style={
            sharp corners,
            colback = black!5, % <-- Modificato qui (5% nero, 95% bianco)
            before skip = 0.2cm,
            after skip = 0.5cm,
            boxrule = .5pt,
            colframe = black,
            % Opzionale: aggiunge un po' di spazio interno
            left=2mm, right=2mm, top=1mm, bottom=1mm
        }
}
\tcolorboxenvironment{formula}{
    formulastyle
}


\renewcommand{\lecture}[2]{%
    \clearpage
    \refstepcounter{lecture} % increase lecture number automatically
    %\refstepcounter{lecture} % Lable referencing

    \section*{\xcharter Mod \thelecture: #1 \\ \small #2}% prints the title unnumbered
    \addcontentsline{toc}{section}{Mod \thelecture: #1 (#2)}% adds it to the table of contents

    \setcounter{section}{0}% reset section numbering
    \setcounter{subsection}{0}% reset subsection numbering
    \setcounter{subsubsection}{0}% reset subsubsection numbering
}